\documentclass[12pt]{article}
\usepackage[utf8]{inputenc}
\usepackage[T1]{fontenc}
\usepackage{hyperref}
\title{Teop Language Resources for K-12 Curriculum Development}
\author{Compiled by Claude Sonnet 5 (claude-sonnet-5) for Andr\'as Kornai}
\date{}
\begin{document}
\maketitle

\section{Introduction}
Teop (ISO 639-3 tio; Glottocode teop1238) is an Oceanic Austronesian language of Bougainville, Papua New Guinea. Speaker estimates range from about 5,000 to 10,000. It is among the better documented languages of the region, thanks to Ulrike Mosel's documentation work.

\section{Primary Educational Resources}
No primers, textbooks or school programmes turned up in search.

\section{Linguistic and Reference Materials}
\begin{itemize}
\item Teop--English dictionary by Ulrike Mosel (Dictionaria, 2019). It has 6,488 entries, many with example sentences and 1,641 with illustrations. It covers the coastal dialect. \url{https://dictionaria.clld.org/contributions/teop}
\item Dictionaria language page for Teop. \url{https://dictionaria.clld.org/languages/teop1238}
\item Teop sketch grammar by Ulrike Mosel with Yvonne Thiesen (2007). The DoBeS archive lists it as an unpublished manuscript. \url{https://www.smg.surrey.ac.uk/media/1147/teop_sketch_grammar_may07.pdf}
\item Overview of Teop in the Northwest Solomonic database (Surrey). \url{https://www.smg.surrey.ac.uk/languages/northwest-solomonic/teop/overview/}
\end{itemize}

\section{Religious and Community Resources}
Nothing found in search.

\section{Digital and Audio Resources}
\begin{itemize}
\item Multi-CAST Teop, a corpus of recorded narratives with audio, at the University of Bamberg. \url{https://fis.uni-bamberg.de/handle/uniba/97606}
\item DoBeS Teop project team page. \url{https://dobes.mpi.nl/projects/teop/team/}
\item DoBeS poster on the Teop project. \url{https://dobes.mpi.nl/posters/print/oceania-teop-print.pdf}
\end{itemize}

\section{Curriculum Development Context}
The dictionary, with its illustrated entries and example sentences, and the narrative corpus are strong sources for a curriculum. No existing school materials were found. Any curriculum should be developed with Teop speakers and the researchers involved.
\end{document}